\section{Methods}
\label{sec:methods}

\subsection{The COSI Hypothesis}

Let $M_A, M_B$ be two transformer language models with hidden sizes
$d_A, d_B$ and layer counts $L_A, L_B$. For a prompt $x_i$ from a probe set
$X = \{x_1, \ldots, x_N\}$, let $a_A^{(\ell)}(x_i) \in \R^{d_A}$ denote the
residual stream activation at layer $\ell$ of $M_A$, pooled over the
sequence dimension to a single vector. Let $X_A^{(\ell)} \in \R^{N \times d_A}$
denote the matrix of pooled activations across all $N$ probes at layer
$\ell$, and analogously for $X_B^{(\ell')}$.

The COSI hypothesis is the following.

\begin{definition}[Approximate invariant subspace]
A $k$-dimensional subspace $V \subseteq \R^d$ is \emph{approximately
$\epsilon$-invariant} under a (possibly nonlinear) map $f: \R^d \to \R^d$ on a
sample $X$ if
$\fnorm{(I - P_V) f(P_V X)} / \fnorm{f(P_V X)} \leq \epsilon$,
where $P_V$ is the orthogonal projection onto $V$.
\end{definition}

\begin{proposition}[COSI hypothesis]
\label{prop:cosi}
For two transformer language models $M_A$ and $M_B$ that have learned
compositional reasoning, there exist layer indices $\ell_A, \ell_B$,
subspace dimension $k$, and approximately $\epsilon$-invariant subspaces
$V_A \subseteq \R^{d_A}$ and $V_B \subseteq \R^{d_B}$ of their respective
forward operators, such that the projected activation matrices
$P_{V_A} X_A^{(\ell_A)}$ and $P_{V_B} X_B^{(\ell_B)}$ admit an orthogonal
alignment with Procrustes residual significantly below the chance level
defined by permutation null.
\end{proposition}

This is a falsifiable empirical claim. The remainder of this section
specifies the protocol that operationalizes it.

\subsection{Activation Extraction}

For each model, we replicate the inner forward pass of the architecture
exactly as implemented in MLX, capturing the residual stream
$h^{(\ell)}(x)$ after each layer block at a set of pre-registered layer
fractions $f \in \{0.25, 0.5, 0.75\}$, mapping to layer indices
$\ell = \max(0, \lfloor f \cdot L \rceil - 1)$. Pooling over the sequence
dimension is performed as last-token (the position immediately preceding
the model's would-be next token), with mean and max pooling reserved as
ablations. Activations are cast from \texttt{bfloat16} to \texttt{float32}
before crossing to host memory. We do not modify the model implementations;
the extraction is a parallel forward pass that uses the same submodules
(embeddings, layers, masks) as the published \texttt{mlx\_lm} forward pass.

\subsection{Subspace Projection}

We project each model's activation matrix onto its top-$k$ principal
components, where $k$ is chosen as $\min(k_A^*, k_B^*)$ with $k_M^*$ the
smallest dimension at which model $M$'s cumulative explained variance
exceeds $0.90$. This ensures both projections retain comparable
information content and that the alignment is performed in matched
dimensions. PCA is computed by truncated SVD on the column-centered
activation matrix.

\subsection{Procrustes Alignment}

In the shared $k$-dimensional subspace, we compute the orthogonal
transformation $R \in O(k)$ that minimizes $\fnorm{X_A - X_B R}$, using
the standard SVD-based solution \cite{schonemann1966generalized}:
$R = U V^\top$, where $U \Sigma V^\top$ is the SVD of $X_B^\top X_A$.
The \emph{Procrustes residual} is
\begin{equation}
\rho(X_A, X_B) = \frac{\fnorm{X_A - X_B R^*}}{\fnorm{X_A}},
\label{eq:procrustes-residual}
\end{equation}
which lies in $[0, \sqrt{2}]$ and equals $0$ for perfect orthogonal
alignment.

\subsection{Null Distributions}

We define two null distributions, both essential.

\textbf{Permutation null.} For each of $S$ random permutations
$\pi: \{1, \ldots, N\} \to \{1, \ldots, N\}$, we recompute the Procrustes
residual on $X_A$ and $X_B^{(\pi)}$, where $X_B^{(\pi)}$ has its rows
permuted according to $\pi$. This breaks the semantic correspondence
between rows while preserving each model's internal representational
geometry. The distribution of residuals over $S$ permutations defines
the chance level for ``no semantic alignment.'' We use $S = 1000$.

\textbf{Random-rotation null.} For each of $S$ random orthogonal matrices
$R_{\text{rand}} \in O(k)$ (sampled uniformly via QR decomposition of a
Gaussian matrix, with sign correction \cite{mezzadri2007generate}), we
compute $\fnorm{X_A - X_B R_{\text{rand}}} / \fnorm{X_A}$. This is a
weaker null---it does not optimize over $R$---and is included for
completeness rather than as the primary test.

\subsection{Significance Criterion}

The pre-registered significance criterion, declared before observing
data, is the conjunction:
\begin{enumerate}
    \item the observed Procrustes residual is below the first percentile
        of the permutation null distribution;
    \item the standardized z-score of the observed residual against the
        permutation null mean is at most $-2$.
\end{enumerate}
A run that fails either condition is reported as a null result.

\subsection{Ablations and Sensitivity Sweeps}

Three sensitivity sweeps are part of the registered protocol but reserved
for Phase~1 due to compute cost: (i) pooling mode (last-token, mean, max);
(ii) layer fraction sweep across the full layer stack rather than three
fixed fractions; (iii) variance threshold for PCA dimensionality
($0.85$, $0.90$, $0.95$). We additionally compute centered kernel alignment
(CKA) \cite{kornblith2019similarity} as a cross-validation metric.
Discrepancies between Procrustes and CKA verdicts will be reported as
informative findings rather than treated as failures of either method.
